% ============================================================
% JMGM Introduction — F3-b re-scope (2026-10-06)
% Central question: do semi-empirical NTO descriptors retain rank information on
% measured TADF gaps that fingerprints lose on publications absent from training?
% Numbers trace to ledger A-006, A-007, A-016--A-020 and the Results section.
% ============================================================

\section{Introduction}\label{sec:introduction}

\subsection{Predicting the singlet--triplet gap of TADF emitters}

Thermally activated delayed fluorescence (TADF) emitters---organic donor--acceptor
molecules that harvest both singlet and triplet excitons for near-unity internal
quantum efficiency~\cite{Uoyama2012}---are the heavy-metal-free route to tunable,
efficient light sources for displays and for human-centric and horticultural
lighting~\cite{Yokoyama_2018, Bugbee2016}. An efficient emitter has to meet several
requirements at once: a small singlet--triplet gap
($\Delta E_\text{ST} < \qty{0.2}{\electronvolt}$), fast reverse intersystem crossing
($k_\text{RISC} > \qty{1e6}{\per\second}$), and emission matched to the target band.
This paper addresses only the first. The rate and device-efficiency requirements need
spin--orbit-coupling and device-level calculations beyond the cheap structural data
used here, and we make no quantitative claim about them (\cref{subsec:limitations}).

\subsection{Machine learning on computed and measured gaps}

Designing such emitters means searching a large space of donor--acceptor
combinations for a small $\Delta E_\text{ST}$. The direct route---time-dependent
DFT on excited states---scales steeply with system size~\cite{Curtarolo2013},
so exhaustive screening of more than a few hundred candidates is impractical.
Machine learning offers a cheaper surrogate~\cite{Butler2018, Pyzer-Knapp2022},
and several groups have predicted TADF gaps from molecular
descriptors~\cite{Gomez-Bombarelli2016, Barneschi2024}. Those low reported errors
are obtained against \emph{computed} (TD-DFT or higher) reference gaps. Nigam
et al.~\cite{Nigam2024} also worked with computed labels, gaps simulated at the
$\omega$B2PLYP$'$ level, and replaced an early regression model with a good/bad
classifier (\cref{subsec:comparison}); their study does not show how models behave
on measured gaps. The closest prior work on measured gaps is that of He
et al.~\cite{He2026}, who trained machine-learning models on \num{1321}
experimental $\Delta E_\text{ST}$ records with measurement-condition features.
Neither study compares a structural fingerprint with a semi-empirical
excited-state representation under molecule, scaffold and source-publication
holdout, which is the comparison made here.
Experimental molecular-property learning is not, by itself, a new approach:
JMGM studies have combined explainable predictors with experimental
ligand--receptor observations~\cite{Ferri2026TAS2R} and modelled experimental
infrared peaks~\cite{Noreen2026IR}. Our question concerns the representation and
the unit that is held out.

One problem recurs in this literature. Accuracy is often quoted against a single random train--test split, which rewards
interpolation within familiar scaffolds and overstates how a model behaves on new
chemistry. Literature-mined corpora sharpen this problem: one publication usually
reports a series of close analogues measured under one protocol, so a random or
scaffold split leaves same-paper siblings of a test molecule in training.

Our earlier work set the stage: Article~1~\cite{Tchapet2024_xTB} benchmarked in JCIM
GFN2-xTB/sTDA against \num{747} experimental molecules
(MAE = \qty{0.162}{\electronvolt}), xTB geometry relaxations, and scaffold
cross-validation, and Article~2~\cite{Tchapet2024_DataDriven}
extracted structure--property design rules from that set. Here we ask: do
semi-empirical natural-transition-orbital (NTO) descriptors retain rank information
on measured TADF singlet--triplet gaps that structural fingerprints lose when
molecules come from publications absent from training? The measured gaps come from
the ChemDataExtractor-generated TADF database of Huang and Cole~\cite{HuangCole2024};
after matching to our descriptor set they cover \num{231} donor--acceptor emitters
from \num{173} publications.

\subsection{Our contributions}

We compare error and rank ordering under three holdout units rather than claim a
new high-accuracy predictor. The answer to the question is yes, within this
corpus: under source-paper holdout the NTO model orders emitters across
publications absent from training with median Spearman $\rho = 0.34$, against
$0.17$ for Morgan fingerprints. Neither representation orders molecules within a
held-out paper, and neither has lower error than the other. The points below
follow the order of the Results.

\paragraph{Fingerprint and NTO models have similar error.}
On \num{231} molecules across \num{212} Bemis--Murcko scaffolds, random forests built
from the 2D structure alone---Morgan fingerprints (MAE = \qty{0.091}{\electronvolt}) and
RDKit descriptors (\qty{0.100}{\electronvolt})---match descriptors derived from a
semi-empirical excited-state calculation (NTO spatial features,
\qty{0.096}{\electronvolt}); the Morgan and NTO models sit only \qtyrange{10}{15}{\percent}
below a mean-value predictor (\qty{0.107}{\electronvolt}). The paired MAE difference is $-\qty{0.0043}{\electronvolt}$ for Morgan minus NTO
(\qty{95}{\percent} CI \qtyrange{-0.0171}{+0.0092}{\electronvolt}). The error
comparison does not move with the holdout unit: under source-paper holdout the
median Morgan-minus-NTO difference is \qty{-0.0004}{\electronvolt}. Unlike a
physics-guided model trained on NTO descriptors against a \emph{computed} reference
gap~\cite{Das2025}, we test the same class of descriptor against the
\emph{measured} gap.

\paragraph{Fingerprint ranking depends on same-publication molecules.}
Over \num{50} random fold assignments, the NTO-minus-Morgan rank difference is
$\Delta\rho = 0.160$ under source-paper holdout (range \numrange{0.07}{0.25}),
against \num{0.073} under scaffold and \num{0.055} under molecule holdout; NTO ranks
higher in \num{19} of \num{20} repeated paper partitions, against \num{11} and
\num{13} of \num{20} for matched scaffold and molecule partitions. The ranges
describe fold-assignment variability, not confidence intervals. NTO $\rho$ stays
between \num{0.337} and \num{0.352} across units; Morgan $\rho$ falls from
\num{0.305} to \num{0.172}. The drop sits on molecules with same-paper siblings:
for them Morgan $\rho$ is \num{0.54} when siblings are in training and \num{0.26}
when the paper is held out, while on the \num{128} molecules from single-molecule
papers it is \num{0.10} under either holdout. NTO is higher on those molecules, but
the seed-averaged NTO-minus-Morgan contrast there ($\Delta\rho = \num{0.16}$,
\qty{95}{\percent} CI \numrange{-0.05}{0.37}) is not resolved. RDKit 2D
descriptors lose rank under paper holdout as Morgan does ($\rho$ \num{0.303} to
\num{0.209}), so among the three sets tested the invariance belongs to the NTO
descriptors. NTO orders emitters across publications absent from training; no
representation orders molecules within a held-out paper. Siblings share close
structure and measurement protocol, and these tests cannot separate the two, so we
report the dependence as an association.

\paragraph{Label scatter leaves room the models do not use.}
Independent experimental reports for the same compound scatter by
\qty{0.072}{\electronvolt} RMS (the largest within-molecule standard deviation is
\qty{0.56}{\electronvolt}). For the \num{102} molecules with several reports, the
median label has an RMS standard error of \qty{0.049}{\electronvolt}; counting the
\num{129} single-report labels at single-report scatter, label noise over all
\num{231} molecules is about \qty{0.072}{\electronvolt} RMS.
Individual discrepancies in reported $\Delta E_\text{ST}$ have been noted for
TADF emitters~\cite{Shuaib2026}; here we quantify inter-report variability
across the \num{231}-molecule corpus drawn from \num{173} papers.
Applying the noise-ceiling framework of Crusius et al.~\cite{Crusius2025} with
per-molecule label noise gives $R^2_\text{max} \approx 0.80$--\num{0.86}
(\cref{subsec:ceilings}), well above the observed $R^2 \approx 0.26$, so label
noise is not the dominant limit and it does not establish a
\qty{0.1}{\electronvolt} MAE floor. Computed references face a separate limit:
B3LYP and CAM-B3LYP gaps differ by up to \qty{0.58}{\electronvolt}. On this corpus
neither a higher-capacity tabular learner (gradient boosting, MLP, SVR) nor the one
frozen pretrained molecular-transformer embedding we tested beats the random forest
(the latter reaches MAE \qtyrange{0.109}{0.113}{\electronvolt}). Low-gap enrichment
on held-out labelled data (\numrange{1.2}{1.4}$\times$ over random) and a filtered
virtual library are reported in the SI; neither supports molecule-level selection,
and we make no device-level performance claims.

This study builds on our earlier work---\textbf{Article~1}~\cite{Tchapet2024_xTB}
(xTB/sTDA validation) and \textbf{Article~2}~\cite{Tchapet2024_DataDriven} (design
rules)---by testing how the holdout unit changes the comparison on measured labels.
The measured labels come from the Huang--Cole TADF database~\cite{HuangCole2024}.
Data, code and the analysis ledger for the three related articles are archived under
Zenodo concept DOI 10.5281/zenodo.17436069 (latest version) and mirrored on GitHub; the
scripts and outputs of the holdout analyses reported here are in that archive and are
also supplied as Supplementary Data (see the Data and code availability statement).

The contribution of this work is a matched comparison of three holdout units
(molecule, scaffold and source publication) on the same \num{231} measured gaps,
which shows that fingerprint and RDKit ranking of measured TADF gaps depends on
same-publication molecules in training while semi-empirical NTO ranking does not.
The error null carries its own information. Because the NTO and fingerprint models
do not differ in error, the sTDA calculation is not
justified by accuracy; its measurable value is ranking across publications. The
\qtyrange{10}{15}{\percent} gain over the mean-value predictor is the margin a new
structure representation has to beat on these labels.
